\ifx\danslelivre\undefined
\documentclass[10pt]{../fichiers-configuration-latex/smfart}
\usepackage{ifthen}
\newboolean{english}
\newboolean{francais}
\setboolean{english}{false}
\setboolean{francais}{true}
\input{../fichiers-configuration-latex/paquets}
\selectlanguage{francais}
\input{../fichiers-configuration-latex/lemme-proposition-theoreme}
\input{../fichiers-configuration-latex/numerotation}
\input{../fichiers-configuration-latex/macros}

\author{Luc Illusie}
\title{XI. Produits orientés}

\begin{document}
\begin{center}
XI. Produits orientés\\
Luc Illusie
\end{center}

\else
\setboolean{english}{false}
\setboolean{francais}{true}
\selectlanguage{francais}
\chapter{Produits orientés}
\addtocontents{toc}{par Luc Illusie}
\begin{center}
Luc Illusie
\end{center}
\fi

On fixe un univers $\mathscr{U}$. Sauf mention du contraire, les sites (resp. topos) considérés
seront des $\mathscr{U}$-sites (resp. topos), et « petit  » signifiera $\mathscr{U}$-petit.

\section{Construction des produits orientés} La construction suivante est due à Deligne
\cite{AG_Tokyo-Kyoto@Laumon} :

\subsection{}\label{11-site} Soient $f : X \ra S$, $g : Y \ra S$ des
morphismes de topos. On suppose que $X$, $Y$, $S$ ont des sites de définition  $C_1$, $C_2$, $D$,
admettant des limites projectives finies, et associés à des prétopologies, et que $f^{*}$, $g^{*}$
prolongent des foncteurs continus entre sites, et commutant aux limites projectives finies. Soit $C$
le site suivant :

(i) $C$ est la catégorie des couples de morphismes $U \ra V \la W$ au-dessus de $X \ra S \la Y$, où
$U \ra V$ (resp. $V \la W$) désigne un morphisme $U \ra f^{*}V$ (resp. $g^{*}V \la W$) de $C_1$
(resp. $C_2$) et $V$ est un objet de $D$.

(ii) $C$ est muni de la topologie définie par la prétopologie engendrée par les familles couvrantes
$(U_i \ra V_i \la W_i) \ra (U \ra V \la W)$ $(i \in I)$ du type suivant :

(a) $V_i = V$, $W_i = W$ pour tout $i$, et la famille $(U_i \ra U)$ est couvrante ;

(b) $U_i = U$, $V_i = V$ pour tout $i$, et la famille $(W_i \ra W)$ est couvrante ;

(c) $(U' \ra V' \la W') \ra (U \ra V \la W)$, où $U' = U$ et $W' \ra W$ est déduit par changement de
base d'un morphisme $V' \ra V$ de $D$. 

On note $\gtilde C$ le topos des faisceaux sur $C$.

\begin{lemme}\label{11-faisceau}
Soit $F$ un préfaisceau sur $C$. Pour que $F$ soit un
faisceau il faut et il suffit que les deux conditions
suivantes soient vérifiées :

(i) pour toute famille couvrante $(Z_i \ra Z)$ de $C$ du type (a) ou (b), la suite $F(Z) \ra
\prod_{i \in I} F(Z_i) \rightrightarrows \prod_{(i,j) \in I \times I} F(Z_i \times_Z Z_j)$ est
exacte ;

(ii) pour toute famille couvrante $(U' \ra V' \la W') \ra (U \ra V \la W)$ du type (c),
l'application $$ F(U \ra V \la W) \ra F(U' \ra V' \la W') $$ est bijective. En particulier, si l'on
note $(-)^a$ le foncteur faisceau associé, pour toute famille couvrante $(Z' \ra Z)$ du type (c), le
morphisme de faisceaux associés $Z'^a \ra Z^a$ est un isomorphisme.
\end{lemme}

La nécessité est triviale pour (i), et pour (ii), il suffit d'observer que le morphisme diagonal $$
(U \ra V' \la W') \ra (U \ra V' \times_V V' \la W' \times_W W') $$ est un morphisme couvrant (du
type (c)), qui égalise la double flèche $$ (U \ra V' \times_V V' \la W' \times_W W')
\rightrightarrows (U \ra V' \la W').  $$ Pour la suffisance, on note que les familles couvrantes
$(Z_i \ra Z)$ de type (a), (b), ou (c) sont stables par changement de base $Z' \ra Z$, et on
applique \sga{4}{ii}{2.3}.

\subsection{}\label{11-tau} Notons $e_X$ (resp. $e_Y$, $e_S$) l'objet final de $C_1$ (resp. $C_2$,
$D$). On a des projections naturelles $$ p_1 : \gtilde C \ra X, p_2 : \gtilde C \ra Y $$ données par
$$ p_1^*(U) = (U \ra e_S \la e_Y), p_2^*(W) = (e_X \ra e_S \la W).  $$ On a par ailleurs un
morphisme canonique $$ \tau : gp_2 \ra fp_1 $$ donné par le morphisme de foncteurs $\tau : (gp_2)_*
\ra (fp_1)_*$ défini de la façon suivante : pour un faisceau $F$ sur $C$, et un objet $V$ de $S$, $$
\tau : ((gp_2)_*F)(V) \ra ((fp_1)_*F)(V) $$ est le composé $$ F(e_X \ra e_S \la g^*V) \ra F(f^*V \ra
V \la g^*V) \ra F(f^*V \ra e_S \la e_Y), $$ où la première flèche est induite par la localisation
$(f^*V \ra V \la g^*V) \ra (e_X \ra e_S \la g^*V)$, et la seconde est l'inverse de l'isomorphisme
donné par \ref{11-faisceau}, relativement au morphisme de type (c) $(f^*V \ra V \la g^*V \ra (f^*V
\ra e_S \la e_Y)$.

\begin{theoreme}\label{11-existence} Soit $T$ un topos muni de morphismes $a : T
\ra X$, $b : T \ra Y$ et d'un morphisme $t : gb \ra fa$. Il existe alors un triplet $(h : T \ra
\gtilde C, \alpha : p_1h \iso a, \beta : p_2h \iso b$, unique à isomorphisme unique près, tel que le
composé $$ \xymatrix{gb \ar[r]^{\beta^{-1}} & gp_2h \ar[r]^{\tau} & fp_1h \ar[r]^{\alpha} & fa} $$
soit égal à $t$.
\end{theoreme}

Nous aurons besoin, pour la démonstration, du lemme suivant :

\begin{lemme}\label{11-lemme existence}
Soit $Z = (U \ra V \la W)$ un objet de $C$. Avec la notation
de \ref{11-faisceau}, le carré suivant est cartésien :
\begin{equation}
\xymatrix{Z^a \ar[d] \ar[r]
& (p_2^*W)^a \ar[d]^v \\ (p_1^*U)^a \ar[r]^u & ((gp_2^*)V)^a}.  \label{11-carré}
\end{equation}
Dans ce carré, $v$ et les flèches issues de $Z^a$ sont les flèches évidentes, et $u$ est la flèche
composée $$ \xymatrix{(p_1^*U)^a \ar[r] & ((fp_1)^*V)^a \ar[r]^{\tau} & ((gp_2^*)V)^a}, $$ où $\tau$
est le composé $$ \xymatrix{(f^*V \ra e_S \la e_Y)^a \ar[r]^{r^{-1}} & (f^*V \ra V \la g^*V)^a
\ar[r]  & (e_X \ra e_S \la g^*V)^a}, $$ $r$ désignant l'isomorphisme $(f^*V \ra V \la g^*V)^a \iso
(f^*V \ra e_S \la e_Y)^a$ de \ref{11-faisceau}.
\end{lemme}

Soit $z : Z \ra Z' = (f^*V \ra V \la
g^*V)$ la projection canonique. Le composé $Z \ra p_1^*U \ra (fp_1)^*V$ se factorise à travers $z$.
Par suite, et par définition de $\tau$, le diagramme $$ \xymatrix{Z^a \ar[r]^z \ar[d] & Z'^a
\ar[d]^{r} \ar[dr] & {}\\ (p_1^*U)^a \ar[r] & ((fp_1^*)V)^a \ar[r]^{\tau} & ((gp_2)^*V)^a } $$ est
commutatif. Comme le composé  $Z \ra p_2^*W \ra (gp_2)^*W$ se factorise aussi à travers $z$, le
carré  \ref{11-carré} est donc commutatif. Celui-ci est le pourtour du diagramme suivant, où les
flèches autres que $\tau$ sont les flèches évidentes : $$ \xymatrix {(f^*U \rightarrow V \rightarrow
g^*W)^a \ar[r]  \ar[d]&(f^*V \rightarrow V \leftarrow g^*W)^a \ar[r] \ar[d]& (e_X \rightarrow e_S
\leftarrow g^*W)^a \ar[d] \\ (f^*U \rightarrow V \rightarrow  g^*V) ^a\ar[r]  \ar[d]&(f^*V
\rightarrow V \leftarrow g^*V)^a \ar[r] \ar[d]^r & (e_X \rightarrow e_S \leftarrow g^*V)^a
\ar[d]^{\Id} \\ (f^*U \rightarrow e_S \leftarrow e_Y)^a \ar[r] &(f^*V \rightarrow e_S \leftarrow
e_Y)^a \ar[r]^{\tau} & (e_X \rightarrow e_S \leftarrow g^*V)^a}.  $$ Chacun des carrés qui le
composent est cartésien. Il en est donc de même de \ref{11-carré}.

\subsection{}\label{11-preuve du th.}
Prouvons \ref{11-existence}. On peut supposer que $a$ et $b$ sont donnés par des morphismes de
sites, la topologie du site de définition de $T$ étant moins fine que la topologie canonique
(d'après \sga{4}{iv}{4.9.4}, on pourrait prendre $T$ lui-même comme site de définition, avec sa
topologie canonique). Par \ref{11-lemme existence} l'unicité est claire : pour $Z = (U \ra V \la W)$
dans $C$, on doit avoir
\begin{equation}
h^*Z = a^*U \times_{(gb)^*V}b^*W, \label{11-déf de h}
\end{equation}
où $a^*U \ra (gb)^*V$ est le composé $\xymatrix{a^*U \ar[r] & (fa)^*V \ar[r]^t &
(gb)^*V}$. Les isomorphismes $\alpha$ et $\beta$ sont alors tautologiques, nous les négligerons dans
le reste de la démonstration. Vérifions que le foncteur $h^*$ donné par \ref{11-déf de h} définit un
morphisme de topos $h$ vérifiant la propriété énoncée en \ref{11-existence}. Comme $h^*$ commute aux
limites projectives finies, pour vérifier que $h^*$ induit un morphisme de topos, il suffit de
vérifier que $h^*$ est continu (\sga{4}{iv}{4.9.1, 4.9.2}). Il est trivial que $h^*$ transforme
familles couvrantes du type (a) ou (b) en familles couvrantes. Par ailleurs, si $(U' \ra V' \la W')
\ra (U \ra V \la W)$ est une famille couvrante du type (c), le carré $$ \xymatrix {b^*W' \ar[r]
\ar[d] &b^*W \ar[d]  \\ (gb)^*V' \ar[r] &(gb)^*V} $$ est cartésien, et par suite $$ a^*U
\times_{(gb)^*V'} b^*W' \rightarrow a^*U \times_{(gb)^*V} b^*W $$ est un isomorphisme. Il reste à
vérifier que $\tau$ induit $t$. Mais par définition, le morphisme de faisceaux défini par
$h^*(\tau)$ appliqué à $V$ est le composé $$ \xymatrix{((fa)^*V)^a \ar[r]^-{r^{-1}} & ((fa)^*V
\times_{(gb)^*V} (gb)^*V)^a \ar[r] & ((gb)^*V)^a}, $$ donc est égal à celui défini par $t$ appliqué
à $(fa)^*V$, ce qui achève la démonstration. 

\begin{definition}
\label{11-produit orienté} Le topos
$\gtilde C$ construit en \ref{11-site} s'appelle le \textit{produit orienté (gauche) de $X$ et $Y$
au-dessus de $S$}, et se note $X \ori_S Y$.  Les morphismes du diagramme $$ \xymatrix{{} & X \ori_S
Y \ar[dl]_{p_1} \ar[dr]^{p_2} & {} \\ X \ar[dr]^f  & {} & Y \ar[dl]_g \\ {} & S & {}} $$ sont reliés
par la 2-flèche $\tau : gp_2 \ra fp_1$. Il découle de \ref{11-existence} que le quadruplet $(X
\ori_S Y, p_1,p_2,\tau)$ est indépendant (à isomorphisme unique près) du choix des sites de
définition $C_1$, $C_2$, $D$.
\end{definition}

On définit de même le produit orienté droit $X \stackrel{\ra}{\times}_S Y$, avec ses projections
canoniques $p_1 : X \stackrel{\ra}{\times}_S Y \ra X$, $p_2 : X \stackrel{\ra}{\times}_S Y \ra Y$ et
la 2-flèche $\tau' : fp_1 \ra gp_2$, qui possède la propriété universelle de \ref{11-existence},
avec $X$ et $Y$ échangés.

\subsection{}\label{11-points} Désignons par $pt$ un topos ponctuel
(catégorie des faisceaux d'ensembles sur un espace réduit à un point). Soient $x : pt \ra X$, $y :
pt \ra Y$ des points de $X$ et $Y$ respectivement, et $u : gy \ra fx$ une 2-flèche. Par
\ref{11-existence}, le triplet $(x,y,u)$ définit un point $z : \mathrm{pt} \ra X \ori_S Y$ tel que
$p_1z \isononcan x$, $p_2z \isononcan y$. Ce point sera noté $(x,y,u)$ (ou parfois $(x,y)$ s'il n'y
a pas de confusion à craindre). Tout point de $X \ori_S Y$ est de cette forme.

\subsection{}\label{11-fonctorialité}
Considérons un diagramme de 1-morphismes de topos
\begin{equation}
\xymatrix{X' \ar[r]^{f'} \ar[d]^u & S' \ar[d]^h & Y' \ar[l]_{g'} \ar[d]^v \\ X
\ar[r]^f & S & Y \ar[l]_{g}} \label{11-fonct1}
\end{equation}
et des 2-flèches $a : hf' \ra fu$, $b
: gv \ra hg'$. Notons $T = X \ori_S Y$, $T' = X' \ori_{S'}Y'$, $p_1 : T \ra X$, $p_2 : T \ra Y$,
$p'_1 : T' \ra X'$, $p'_2 : T' \ra Y'$ les projections canoniques, $\tau : gp_2 \ra fp_1$, $\tau' :
g'p'_2 \ra f'p'_1$ les 2-flèches canoniques. Considérons la 2-flèche composée $$ c :
\xymatrix{gvp'_2 \ar[r]^b & hg'p'_2 \ar[r]^{\tau'} &hf'p'_1 \ar[r]^a &fup'_1}.  $$ D'après
\ref{11-existence}, $c$ définit un diagramme de 1-morphismes
\begin{equation}
\xymatrix{X' \ar[d]^u
&T' \ar[d]^t \ar[l]_{p'_1} \ar[r]^{p'_2} & Y' \ar[d]^v \\ X & T \ar[l]_{p_1}  \ar[r]^{p_2} & Y}
\label{11-fonct2}
\end{equation}
et des 2-isomorphismes $\alpha : p_1t \iso up'_1$, $\beta : p_2t
\iso vp'_2$ rendant commutatif le carré
\begin{equation}
\xymatrix{gp_2t \ar[r]^{\tau}
\ar[d]^{\beta} &  fp_1t \ar[d]^{\alpha} \\ gvp'_2 \ar[r]^{a\tau'b} &fup'_1}.  \label{11-fonct3}
\end{equation}
On dit que le triplet $(t,\alpha,\beta)$ (ou simplement $t : T' \ra T$) de
\ref{11-fonct2} est déduit de \ref{11-fonct1} par \textit{fonctorialité}. On notera $$ t = u \ori_h
v.  $$ On a une compatibilité évidente pour un composé de deux données \ref{11-fonct1}.

\subsection{}\label{11-exemples} Voici quelques exemples.

(a) Dans la situation de \ref{11-existence}, le triplet $(a,b,t)$ définit un diagramme de type
\ref{11-fonct1} $$ \xymatrix{T \ar[d]^a \ar[r]^{\Id} & T \ar[d]^{fa} & T \ar[l]_{\Id} \ar[d]^b \\ X
\ar[r]^f & S & Y \ar[l]_{g}}, $$ avec $t : gb \ra fa$, d'où un morphisme $$ a \ori_{fa} b : T \ori_T
T \ra X \ori_S T.  $$ Par ailleurs, d'après \ref{11-existence}, les flèches identiques de $T$
définissent un morphisme canonique, dit \textit{diagonal} $$ \Delta : T \ra T \ori_T T.  $$ La
1-flèche $h$ de \ref{11-existence} est la composée $$ h = (a \ori_{fa} b)\Delta : T \ra X \ori_S Y.
$$ En particulier, prenant pour $T$ un topos ponctuel, de sorte que $\Delta$ est un isomorphisme, on
a, avec les notations de \ref{11-points} $$ (x,y,u) = x \ori_{fx} y : pt \ra X \ori_S Y.  $$

(b) Dans la situation de \ref{11-produit orienté}, soient $X'$, $S'$, $Y'$ des objets de $X$, $S$,
$Y$ respectivement, et $f' : X' \ra S'$, (resp. $g' : Y' \ra S'$) une flèche au-dessus de $f$ (resp.
$g$), i. e. une flèche $f' : X' \ra f^*(S')$, (resp. $g' : Y' \ra g^*(Y')$). Notons $X' \ori_{S'}
Y'$ l'objet $(X' \ra S' \la Y') = p_1^*(X') \times_{(gp_2)^*(S')} p_2^*Y'$ de $X \ori_S Y$,
\textit{cf.} \ref{11-lemme existence}, On en déduit un diagramme 2-commutatif de1- flèches
naturelles

\begin{equation}
\xymatrix{{} & (X \ori_S Y)_{/(X' \ori_{S'}Y')} \ar[dl] \ar[dr] & {} \\
X_{/X'} \ar[r] \ar[d] & S_{/S'} \ar[d] & Y_{/Y'} \ar[l] \ar[d] \\ X \ar[r] & S & Y \ar[l]},
\label{11-localisation1}
\end{equation}
où la notation $(-)_{/-}$ désigne un topos localisé.
D'après \ref{11-existence}, la partie supérieure de \ref{11-localisation1} définit un 1-morphisme
\begin{equation} m :  (X \ori_S Y)_{/(X' \ori_{S'}Y')} \ra X_{/X'} \ori_{S_{/S'}} Y_{/Y'}.
\label{11-localisation2}
\end{equation}
Il résulte de \ref{11-déf de h} que $m$ est une équivalence,
par laquelle, dans la suite, nous identifierons les deux membres. D'autre part, les carrés
2-commutatifs de \ref{11-localisation1} définissent, d'après \ref{11-fonctorialité}, une flèche de
fonctorialité $$ X_{/X'} \ori_{S_{/S'}} Y_{/Y'} \ra X \ori_S Y, $$ Celle-ci, ou son composé avec
$m$,
\begin{equation}
(X \ori_S Y)_{/(X' \ori_{S'}Y')} \ra X \ori_S Y \label{11-localisation3}
\end{equation}
s'appelle \textit{flèche de localisation}.

\begin{proposition}\label{11-localisation4} Supposons que $X'$ soit l'objet final de $X$ et que $g'$
soit cartésien au-dessus de $g$, i. e. $g' : Y' \iso g^*(S')$. Alors la flèche
\ref{11-localisation3} est une équivalence. 
\end{proposition}

En effet, avec les notations de
\ref{11-localisation1}, il résulte de \ref{11-faisceau} que la flèche $e_X \ori_{S'} Y' \ra e_X
\ori_{e_S} e_Y$ de $X \ori_{S} Y$ est un isomorphisme.

\subsection{}\label{11-tubes} Considérons en
particulier le cas où $S = Y$ est un schéma muni de la topologie étale, $g = \Id$ et $f : X \ra S =
Y$ est l'inclusion d'un sous-schéma fermé de $Y$. Le topos $T = X \ori_Y Y$ joue le rôle d'un
\textit{voisinage tubulaire étale} de $Y$ dans $X$. Les points de $T$ sont les triplets $(x,y,t)$,
où $x$ (resp. $y$) est un point géométrique de $X$ (resp. $Y$) et $t : y \ra x$ une flèche de
spécialisation (\textit{cf.  (\sga{4}{viii}{7.9})}. En d'autres termes, $(x,y,t)$ est la donnée d'un
point géométrique $x$ de $X$, d'une générisation $y_0$ du point fermé (noté encore par abus $x$) du
localisé strict  $X_{(x)}$ de $X$ en $x$ et d'un point géométrique de  $X_{(x)}$ localisé en $y_0$,
ou encore d'une extension séparablement close $y \ra y_0$ du point générique de
$\overline{\{y_0\}}$. Par ailleurs, si $v : Y' \ra Y$ est un voisinage étale de $X$ dans $Y$, i. e.
un diagramme commutatif $$ \xymatrix{{} & Y' \ar[d]^v \\ X \ar[ur] \ar[r] & Y}, $$ où $v$ est étale,
alors, d'après \ref{11-localisation4} le morphisme canonique $$ X \ori_{Y'} Y' \ra X \ori_Y Y $$ est
une équivalence. Ainsi, $T$ ne dépend que du hensélisé de $X$ le long de $Y$ (lorsque celui-ci est
défini, en particulier, pour $Y$ affine, \textit{cf.} \cite{Anneaux@Raynaud}). Nous verrons au
numéro suivant et dans l'exposé XII d'autres propriétés de $T$ précisant cette analogie avec un
voisinage tubulaire.  \section{Tubes et changement de base}

\subsection{}\label{11-topos locaux}
Soit $(S,s)$ un topos ponctué, i. e. un couple formé d'un topos $S$ et d'un point $s : pt \ra S$ de
$S$. Si $(S,s)$ et $(T,t)$ sont des topos ponctués, un morphisme (ponctué) de $(S,s)$ dans $(T,t)$
est un couple $(f,a)$ d'un morphisme $f : S \ra T$ et d'une 2-flèche $a : fs \ra t$. Une 2-flèche $c
: (f,a) \ra (g,b)$ est une 2-flèche $c : f \ra g$ telle que $b(cs) = a$. Si $(S,s)$ est un topos
ponctué, on note $F \mapsto F_s = s^*F$ le foncteur fibre en $s$.

Rappelons qu'un topos ponctué $(S,s)$ est dit \textit{local de centre} $s$ (\sga{4}{vi}{8.4.6}) si,
pour tout objet $F$ de $S$, la flèche naturelle $\Gamma(S,F) \ra F_s$ est bijective. Un morphisme
$(f,a) : (S,s) \ra (T,t)$ de topos ponctués est dit \textit{local} si la 2-flèche $a : fs \ra t$ est
un isomorphisme.

\subsection{}\label{11-section canonique} La construction qui suit est due à
Gabber. Soit $(S,s)$ un topos local de centre $s$. Notons $\varepsilon : S \ra pt$ la projection.
Par définition la flèche canonique $\varepsilon_* \ra s^*$ est un isomorphisme. On en déduit un
isomorphisme
\begin{equation}
\varepsilon^*\varepsilon_* \iso (s\varepsilon)^*.  \label{11-can1}
\end{equation} La flèche d'adjonction $\varepsilon^*\varepsilon_* \ra \Id$ s'identifie donc, par
\ref{11-can1}, à un morphisme $(s\varepsilon)^* \ra \Id$, i. e. à une 2-flèche
\begin{equation}
c_s :
\Id  \ra s\varepsilon \label{11-can2}
\end{equation}
entre les 1-morphismes $\Id : S \ra S$ et
$s\varepsilon : S \ra S$. Si $F$ est un objet de $S$, $(s\varepsilon)^*F$ est le faisceau constant
sur $S$ de valeur $F_s = \varepsilon_*F = \Gamma(S,F)$. Si $U$ est un objet connexe de $S$, le
morphisme $(s\varepsilon)^*F \ra F$ induit sur $\Gamma(U,-)$ le morphisme de restriction
$\Gamma(S,F) \ra \Gamma(U,F)$. Le composé $c_ss : s \ra s$ est l'identité : $(s\varepsilon)^*F \ra
F$ induit l'identité sur les fibres en $s$.

Soient $f : X \ra S$, $g : Y \ra S$ des morphismes de topos, $x : pt \ra X$ un point de $X$, $s = fx
: pt \ra S$ son image dans $S$. Le diagramme
\begin{equation}
\xymatrix{pt \ar[r]^{\Id} \ar[d]^x & pt
\ar[d]^{s} & Y \ar[l]_{\varepsilon g} \ar[d]^{\Id} \\ X \ar[r]^f & S & Y \ar[l]_g}, \label{11-can3}
\end{equation} (où le carré de gauche est 2-commutatif) et la 2-flèche
\begin{equation}
c_sg : g \ra
s\varepsilon g \label{11-can4}
\end{equation}
sont une donnée de type \ref{11-fonct1}. Pour un objet
$F$ de $S$, $(s\varepsilon g)^*F$ est le faisceau constant sur $Y$ de valeur $F_s = \Gamma(S,F)$, et
la flèche $(s\varepsilon g)^*F \ra g^*F$ est la composée $\Gamma(S,F)_Y \ra \Gamma(Y,g^*F)_y \ra
g^*F$. Notons que, par \ref{11-existence}, le produit orienté $pt \ori_{pt} Y$ s'identifie
canoniquement à $Y$, avec $p_1 = \Id : Y \ra Y$. De \ref{11-can3} et \ref{11-can4} on déduit donc un
diagramme de type \ref{11-fonct2} :
\begin{equation}
\xymatrix{pt \ar[d]^x & Y \ar[d]^{\sigma}
\ar[l]_{\varepsilon g} \ar[r]^{\Id} & Y \ar[d]^{\Id} \\ X & X \ori_S Y \ar[l]_{p_1} \ar[r]^{p_2} & Y},
\label{11-can5}
\end{equation}
en d'autres termes, une section $\sigma : Y \ra X \ori_S Y$ de $p_2$
telle que $p_1\sigma = x\varepsilon g$. On peut voir cette section comme étant définie, via
\ref{11-existence} par le couple de morphismes $x\varepsilon g : Y \ra X$, $\Id : Y \ra Y$ et la
2-flèche $c_sg : g \ra fx\varepsilon g = s\varepsilon g$. On dit que $\sigma$ est la \textit{section
canonique} définie par le point $x$. Par composition avec $p_2{}_*$, la flèche d'adjonction $\Id \ra
\sigma_*\sigma^*$ donne une flèche canonique
\begin{equation}
\gamma : p_2{}_* \ra \sigma^*.
\label{11-can6}
\end{equation}
Le résultat suivant est dû à Gabber :

\begin{proposition}\label{11-can7} Soient $f : (X,x) \ra (S,s)$ un morphisme local de topos locaux
($fx = s$), et $g : Y \ra S$ un morphisme de topos. Soit $y : pt \ra Y$ un point de $Y$. Pour tout
objet $F$ de $X \ori_S Y$, $\gamma$ (\ref{11-can6}) induit un isomorphisme $$ \gamma_y :
(p_2{}_*F)_y \iso (\sigma^*F)_y.  $$
\end{proposition}

Soit $t = \sigma y : pt \ra T $ le point de
$T = X \ori_S Y$ image de $y$ par $\sigma$. Ce point est défini (\textit{cf.} 1.4) par le triplet
$(x,y,u)$, où $u : gy \ra fx = fx\varepsilon gy = s\varepsilon g y = s$ est déduit de \ref{11-can4}.
On a $(\sigma^*F)_y = F_t$. Par définition, $$ F_t = \mathrm{colim}_z F(U \ra V \la W), $$ où $z :
pt \ra (U \ra V \la W)$ parcourt les voisinages de $t$ dans $T$. Comme $X$ et $S$ sont locaux, les
voisinages de $t$ de la forme $\sigma w : pt \ra (e_X \ra e_S \la W)$, où $w : pt \ra W$ est un
voisinage de $y$ dans $Y$ forment un système cofinal. Donc $$ F_t = \mathrm{colim}_w F(e_X \ra e_S
\la W), $$ où $z = \sigma w : pt \ra (e_X \ra e_S \la W)$ parcourt les voisinages précédents, avec
$U = e_X$, $V = e_S$.  Par ailleurs, $$ (p_2{}_*F)_y = \mathrm{colim}_w F(e_X \ra e_S \la W), $$ où
$w : pt \ra W$ parcourt les voisinages de $y$ dans $Y$. La flèche $\gamma_y$ est la restriction
naturelle. C'est donc un isomorphisme. 

\begin{corollaire2}
\label{11-can8} Sous les hypothèses de
\ref{11-can7}, si $S$ a assez de points, en particulier, si $S$ est localement cohérent
(\sga{4}{vi}{9.0}), $\gamma$ (\ref{11-can6}) est un isomorphisme.
\end{corollaire2}

Il est plausible que l'hypothèse d'existence d'assez de points soit superflue. Celle-ci sera
cependant satisfaite dans les applications que nous avons en vue.

\begin{corollaire2}\label{11-can9} Sous les hypothèses de \ref{11-can7}, supposons $Y$ local de
centre $y$. Alors $X \ori_S Y$ est local de centre $\sigma(y)$. 
\end{corollaire2}

En effet, on a
alors $(p_2{}_*F)_y = \Gamma(Y,p_2{}_*F) = \Gamma(X \ori_S Y,F)$, et $\gamma_y$ s'identifie à la
restriction $ \Gamma(X \ori_S Y,F) \ra F_{\sigma(y)}$.

Notons que, si $f : X \ra S$ est un morphisme local de schémas strictement locaux, $g : Y \ra S$ un
morphisme de schémas strictement locaux, le produit fibré schématique $X \times_S Y$ n'est pas en
général strictement local, ni même local, même si $g$ est local.

Le résultat ci-après est dû également à Gabber :

\begin{theoreme}
\label{11-ch. base} Soient $f : X
\ra S$, $g : Y \ra S$ des morphismes de topos, $T = X \ori_S Y$, $p_1 : T \ra X$, $p_2 : T \ra Y$
les projections canoniques, $\tau = gp_2 \ra fp_1$ la 2-flèche canonique. On suppose $X$, $Y$, $S$
cohérents et $f$, $g$ cohérents (\sga{4}{vi}{2.3,2.4.5,3.1}). Soit $\Lambda$ un anneau. Alors, pour
tout $F \in D^+(Y,\Lambda)$, la flèche de changement de base, déduite de $\tau$,
\begin{equation}
f^*Rg_*F \ra Rp_1{}_*p_2^*F \label{11-ch. base1}
\end{equation}
est un isomorphisme (de
$D^+(X,\Lambda)$).
\end{theoreme}

Nous aurons besoin du lemme suivant, qui généralise \cite[9.1]{Cycles@Orgogozo} :

\begin{lemme}\label{11-cohérent} Sous les hypothèses de \ref{11-ch. base}, $T$ est cohérent, et les
projections $p_1$, $p_2$ sont des morphismes cohérents.
\end{lemme}

Le topos $X$ (resp. $Y$, resp. $S$) admet une (petite) famille génératrice $C_1$ (resp. $C_2$, resp.
$D$) formée d'objets cohérents, stable par limites projectives finies (\sga4{vi}{2.4.5}). Comme $f$
et $g$ sont cohérents, $f^*V$ (resp. $g^*V$) est cohérent si $V$ est dans $D$ (\sga4{vi}{3.2}). La
sous-catégorie pleine de $X$ (resp. $Y$, resp.  $S$ correspondante, munie de la topologie induite,
est un site de définition de $X$ (resp. $Y$, resp. $S$). Soit $C$ la catégorie définie comme en
\ref{11-site}, munie de la topologie définie par la prétopologie engendrée par les familles finies
de type (a) et (b) et les familles de type (c). Elle est stable par limites projectives finies, et
est un site de définition de $T$. Il suffit donc de montrer que tout objet de $C$ est quasi-compact
(\sga4{vi}{2.4.5}). Pour cela, notons $\mathscr{P}$ la prétopologie définie en \ref{11-site}.
Décrivons $\mathscr{P}$. Pour chaque objet $Z = (U \ra V \la W)$ de $C$, notons $\mathrm{Cov}(Z)$
l'ensemble des familles $(Z_i \ra Z)_{i \in I}$ obtenues par composition d'un nombre fini de
familles de type (c) et de familles (finies) de type (a) et (b). En particulier, l'ensemble $I$ est
fini. Par définition, la donnée des $\mathrm{Cov}(Z)$ vérifie les axiomes PT0, PT2 et PT3 de
\sga4{ii}{1.3}.  L'axiome PT1 (stabilité par changement de base) est également vérifié, les familles
de type (c), ainsi que les familles finies de type (a) (resp. (b)) étant stables par changement de
base, et le changement de base commutant à la composition des familles. La donnée des
$\mathrm{Cov}(Z)$ est donc une prétopologie, et par définition, c'est la prétopologie $\mathscr{P}$.
Comme les familles appartenant à $\mathrm{Cov}(Z)$ sont finies, tout objet de $C$ est
automatiquement quasi-compact, comme annoncé. La cohérence des projections $p_1$ et $p_2$ en
découle.


\begin{remarque}\label{11-g suffit} Gabber sait montrer que la conclusion de \ref{11-cohérent} vaut
sous les seules hypothèses que $X$, $Y$, $S$ et $g$ sont cohérents. Nous n'aurons pas besoin de
cette généralisation.
\end{remarque}

\subsection{}\label{11-preuve de ch. base} Prouvons
\ref{11-ch. base}. Comme $X$ est cohérent, donc possède assez de points, il suffit de vérifier que,
pour tout point $x : pt \ra X$ de $X$, la fibre en $x$ de \ref{11-ch. base1}
\begin{equation}
(f^*Rg_*F)_x \ra (Rp_1{}_*p_2^*F)_x \label{11-(*)}
\end{equation}
est un isomorphisme. Soit $s : pt
\ra S$ l'image de $x$ par $f$. Soit $X_{(x)}$ (resp. $S_{(s)}$) le localisé de $X$ (resp. $S$) en
$x$ (resp. $s$). Rappelons (\sga{4}{vi}{8.4.2}) que $X_{(x)}$ (resp. $Y_{(y)}$) (noté
$\mathrm{Loc}_x(x)$ (resp. $\mathrm{Loc}_s(S)$) dans \textit{loc. cit.}) est la limite projective $$
X_{(x)} = \mathrm{Limtop}_{U \in Vois(x)} X_{/U}, $$ (resp.  $$ S_{(s)} = \mathrm{Limtop}_{V \in
Vois(s)} S_{/V}), $$ où $U$ (resp. $V$) parcourt la catégorie cofiltrante $Vois(x)$ (resp. $Vois(s))
$des voisinages de $x$ (resp. $s$) dans $X$ (resp. $S$). Comme $X$ (resp. $S$) est cohérent, on peut
d'ailleurs se borner aux $U$ (resp. $V$) qui sont cohérents, les morphismes de transition étant
alors automatiquement cohérents. C'est ce que nous ferons dans la suite, notant encore $Vois(x)$
(resp. $Vois(s)$) la sous-catégorie pleine formée des $U$ (resp. $V$) cohérents. Le topos $X_{(x)}$
(resp. $S_{(s)}$) est un topos local, au-dessus de $X$ (resp. $S)$, dont l'image du centre est $x$
(resp. $s$) (\sga{4}{vi}{8.4.6}). Le morphisme $f$ induit un morphisme local $f_{(x)} : X_{(x)} \ra
S_{(s)}$. Définissons de même $$ T_{(x)} = \mathrm{Limtop}_{U \in Vois(x)} T_{/p_1^*U}, $$ $$
Y_{(s)} = \mathrm{Limtop}_{V \in Vois(s)} Y_{/g^*V}), $$ de sorte qu'on obtient un carré
\begin{equation} \xymatrix{T_{(x)} \ar[d]^{p_1} \ar[r]^{p_2} & Y_{(s)} \ar[d]^g \\ X_{(x)} \ar[r]^f
& S_{(s)}}, \label{11-carré localisé}
\end{equation}
avec une 2-flèche $\tau : gp_2 \ra fp_1$. Par
la compatibilité de la formation des produits orientés à la localisation (\ref{11-localisation2}),
la flèche
\begin{equation}
T_{(x)} \ra X_{(x)} \ori_{S_{(s)}} Y_{(s)}, \label{11-(0)}
\end{equation}
déduite de ce carré par \ref{11-existence} est une équivalence. D'après \ref{11-cohérent}, les topos
$X_{/U}$, $S_{/V}$, $T_{/p_1^*U}$, $Y_{/g^*V}$ sont cohérents, les flèches de transition des
systèmes projectifs $X_{/U}$, $S_{/V}$ sont cohérents, et les morphismes $g : Y_{/g^*V} \ra S_{/V}$,
$p_1 : T_{/p_1^*U} \ra X_{/U}$ sont cohérents. On est donc dans les conditions d'application de
\sga4{vi}{8.7.3}, qui, compte tenu de ce que $S_{(s)}$ et $X_{(x)}$ sont locaux, implique que les
flèches canoniques
\begin{equation}
(Rg_*F)_s \ra  R\Gamma(S_{(s)},Rg_*F) \ra R\Gamma(Y_{(s)},F),
\label{11-(1)} \end{equation}
\begin{equation}
(Rp_1{}_*p_2^*F)_x \ra
R\Gamma(X_{(x)},Rp_1{}_*p_2^*F) \ra R\Gamma(T_{(x)},p_2^*F) \label{11-(2)}
\end{equation}
sont des
isomorphismes. Avec les identifications \ref{11-(0)}, \ref{11-(1)} et \ref{11-(2)}, la flèche
\ref{11-(*)} s'identifie à la fibre en $x$ de la flèche de changement de base (déduite de $\tau$) du
carré \ref{11-carré localisé}. Cette flèche s'écrit
\begin{equation}
R\Gamma(Y_{(s)},F)  \ra
R\Gamma(T_{(x)},p_2^*F).  \label{11-(3)}
\end{equation}
On a : 

\medskip (*) : \textit{La flèche \ref{11-(3)} s'identifie canoniquement à la flèche de fonctorialité
définie par $p_2$}.  \medskip

Pour le vérifier, on peut supposer $S$ et $X$ locaux, de centres respectifs $s$ et $x$, et $f$
local. Par définition, \ref{11-(3)} est la flèche composée $$ \xymatrix{R\Gamma(S,Rg_*F)
\ar[r]^{\ref{11-(3)}} \ar[d]^{\alpha} & R\Gamma(T,p_2^*F) \\ R\Gamma(S,Rg_*Rp_2{}_*p_2^*F)
\ar[r]^{R\Gamma(S,\tau)} & R\Gamma(S,Rf_*Rp_1{}_*p_2^*F) \ar[u]_{\beta}}, $$ où la flèche $\alpha$
est définie par la flèche d'adjonction $adj : F \ra Rp_2{}_*p_2^*F$ (et donc est la flèche de
fonctorialité $R\Gamma(Y,F)  \ra R\Gamma(T,p_2^*F)$ définie par $p_2$), et $\beta$ est
l'isomorphisme canonique de transitivité relatif à $fp_1 : T \ra S$. Or, par définition de $\tau$
(\ref{11-tau}), pour tout faisceau $G$ sur $T$, la flèche $$ \Gamma(S,\tau) : \Gamma(S,(gp_2)_*G)
\ra \Gamma(S,(fp_1)_*G) $$ est l'identité. Il en est donc de même de la flèche horizontale
inférieure du diagramme ci-dessus, ce qui prouve (*). Il reste à prouver que \ref{11-(3)} est un
isomorphisme. En fait, la flèche
\begin{equation}
adj : F \ra Rp_2{}_*p_2^*F \label{11-adj}
\end{equation} est un isomorphisme. Pour le voir, il suffit d'observer que, compte tenu de la
description de $\gamma$ donnée en \ref{11-can9}, le composé $$ \xymatrix{F \ar[r]^{adj}  &
Rp_2{}_*p_2^*F \ar[r]^{\gamma} & \sigma^*p_2^*F = F}, $$ où $\sigma$ est la section de $p_2$ définie
en \ref{11-can5} et $\gamma$ l'isomorphisme de \ref{11-can8}, est l'identité. Ceci achève la
démonstration de \ref{11-ch. base}.

\begin{remarques}\label{11-remarques}
\renewcommand{\theenumi}{(\arabic{enumi})}
\begin{enumerate}
\item Supposons que les données de \ref{11-ch. base} proviennent de morphismes de schémas, munis de la
topologie étale, avec $X$, $S$, $Y$ cohérents et $f$ et $g$ cohérents. Les points $x$, $s$ sont des
points géométriques, et les localisés $X_{(x)}$, $S_{(s)}$ des localisés stricts. Si $f$ est une
immersion fermée, la flèche d'adjonction \ref{11-adj} est un isomorphisme (on peut en effet supposer
$S$ strictement local, et il en et alors de même de $X$). Supposons de plus que $Y = S$, $g = \Id_S$
comme en \ref{11-tubes}. On a vu en \textit{loc. cit.} que $T = X \ori_S S$ joue le rôle d'un
voisinage tubulaire de $X$ dans $S$. Soient $j : S^*  = S - X \ra S$ l'ouvert complémentaire de $X$,
et $T^* = X \ori_S S^* = T_{/(e_X \ra e_S \la S^*)}$ le topos induit. Alors  $T^* $ joue le rôle
d'un \textit{voisinage tubulaire épointé} de $X$ dans $S$ : pour $F \in D^+(S^*,\Lambda)$, on a, par
\ref{11-ch. base}, $$ f^*Rj_*F \iso Rp_1{}_*p_2^*F.  $$

\item Sans l'hypothèse de cohérence sur $g$, la conclusion de \ref{11-ch. base} peut être en
défaut, comme le montre l'exemple suivant, dû à Gabber. Soient $X$ un espace topologique connexe,
non vide, $i : Y \ra X$ l'inclusion d'un fermé non vide distinct de $X$, $j : U= X - Y \ra X$
l'inclusion de l'ouvert complémentaire. Alors $i^*j_*\ZZ$ est non nul. Mais, si tout point de $U$ a
un voisinage dont l'adhérence dans $X$ ne rencontre pas $Y$, alors le produit orienté $Y \ori_X U$
est vide. C'est le cas par exemple, si $X$ est le segment $[0,1]$ et $Y$ le point $\{0\}$. 

\item Sous les hypothèses de \ref{11-ch. base}, on montre de manière analogue que :
\begin{itemize}

\item[(a)] Pour tout faisceau d'ensembles $F$ sur $Y$, la flèche de changement de base $$ f^*g_*F \ra
p_1{}_*p_2^*F $$ est un isomorphisme.

On peut espérer des variantes non abéliennes supérieures :

\item[(b)] Pour tout faisceau en groupes $F$ sur $Y$, la flèche de changement de base $$ f^*R^1g_*F \ra
R^1p_1{}_*p_2^*F, $$ est un isomorphisme (de faisceaux d'ensembles pointés).

\item[(c)] Plus généralement, pour tout champ $F$ sur $Y$, la flèche de changement de base $$ f^*g_*F \ra
p_1{}_*p_2^*F $$ est une équivalence.

\end{itemize}

La vérification de (b) et (c) semble requérir, outre les techniques de réduction des champs aux
gerbes de \cite[III 2.1.5]{CNA@Giraud}, des résultats de passage à la limite pour la cohomologie non
abélienne analogues à ceux de \sga4{vi}{8.7}, pour lesquels nous ne connaissons pas de référence.

\item Gabber sait démontrer la généralisation suivante de \ref{11-ch. base}. Soient $f : X
\ra S$, $g : Y \ra S$ des morphismes de topos, $T = X \ori_S Y$. On suppose que $Y$ et $S$ sont
localement cohérents, et que, pour tout objet cohérent algébrique $V$ de $S$, $g^*V$ est cohérent
algébrique (\sga4{vi}{2.1,2.3}). Alors, pour tout $F \in D^+(Y,\Lambda)$, la flèche de changement de
base \ref{11-ch. base1} est un isomorphisme, et on devrait avoir des résultats analogues dans le cas
non abélien, comme en (3) (a), (b), (c) ci-dessus. Gabber déduit ces résultats d'un théorème général
de changement de base pour certains topos fibrés.
\end{enumerate}
\end{remarques}

\section{Produits fibrés} Les compléments donnés
dans ce numéro et le suivant ne seront pas utilisés dans le reste du volume.

\subsection{}\label{11-site produit fibré} Les produits fibrés de topos ont été construits par
Giraud \cite[3.4]{classtop@Giraud}. La construction suivante est due à Gabber. Soient $f : X \ra S$,
$g : Y \ra S$ des morphismes de topos comme en \ref{11-site}. Soit $D$ le site suivant :

\renewcommand{\theenumi}{(\roman{enumi})}
\begin{enumerate}
\item La catégorie sous-jacente à $D$ est la catégorie $C$ considérée en \ref{11-site} (i). 

\item $D$ est munie de la topologie définie par la prétopologie engendrée par les familles couvrantes
$(U_i \ra V_i \la W_i) \ra (U \ra V \la W)$ $(i \in I)$ de la forme (a), (b), (c) de \ref{11-site}
(ii) et de la forme

(d) $(U' \ra V' \la W') \ra (U \ra V \la W)$, où $W' = W$ et $U' \ra U$ est déduit par changement de
base d'un morphisme $V' \ra V$ du site de définition de $S$.
\end{enumerate}

En d'autres termes, la topologie sur $D$ est la borne supérieure des topologies sur $C$ définissant
les produits orientés $X \ori_S Y$ et $X \stackrel{\ra}{\times}_S Y$. 

D'après \ref{11-faisceau}, pour qu'un préfaisceau $F$ sur $D$ soit un faisceau, il faut et il suffit
que $F$ vérifie les conditions d'exactitude habituelles relatives aux familles couvrantes de type
(a) et (b), et que, pour toute famille couvrante $Z' \ra Z$ de type (c) ou (d), $F(Z) \ra F(Z')$
soit un isomorphisme. 

Soit $\gtilde D$ le topos des faisceaux sur $D$. On a des projections naturelles $$ p_1 : \gtilde D
\ra X, p_2 : \gtilde D \ra Y $$ données par les mêmes formules qu'en \ref{11-tau}, et la construction
de $\tau$ en (\textit{loc. cit.}) donne un \textit{isomorphisme}
\begin{equation}
\varepsilon : gp_2
\iso fp_1.  \label{11-epsilon} \end{equation}

\begin{theoreme}\label{11-existence produit fibré}
Soit $T$ un topos muni de morphismes $a : T \ra X$, $b : T \ra Y$ et d'un isomorphisme $t : gb \iso
fa$. Il existe alors un triplet $(h : T \ra \gtilde D, \alpha : p_1h \iso a, \beta : p_2h \iso b$,
unique à isomorphisme unique près, tel que le composé $$ \xymatrix{gb \ar[r]^{\beta^{-1}} & gp_2h
\ar[r]^{\varepsilon} & fp_1h \ar[r]^{\alpha} & fa} $$ soit égal à $t$.
\end{theoreme}

La démonstration est analogue à celle de \ref{11-existence}. Le foncteur $h^*$ est encore donné par la
formule \ref{11-déf de h}. Comme $t$ est un isomorphisme, on a $$ a^*U \times_{(gb)^*V}b^*W = a^*U
\times_{(fa)^*V} b^*W, $$ où $b^*V \ra (fa)^*V$ est le composé $\xymatrix{b^*V \ar[r] & (gb)^*V
\ar[r]^{t^{-1}} & (fa)^*V}$. Il s'ensuit que $h^*$ transforme famille couvrante de type (d) en
famille couvrante, et l'on conclut comme dans \ref{11-preuve du th.}.

\begin{definition}\label{11-produit fibré} Le topos $\gtilde D$ s'appelle \textit{produit fibré de
$X$ et $Y$ au-dessus de $S$}, et se note $X \times_S Y$. Les morphismes du diagramme $$ \xymatrix{{}
& X \times_S Y \ar[dl]_{p_1} \ar[dr]^{p_2} & {} \\ X \ar[dr]^f  & {} & Y \ar[dl]_g \\ {} & S & {}}
$$ sont reliés par le 2-isomorphisme $\varepsilon : gp_2 \ra fp_1$.
\end{definition}

\begin{exemples}\label{11-exemples produits fibrés}
\renewcommand{\theenumi}{(\arabic{enumi})}
\begin{enumerate}
\item \textit{Espaces topologiques}. Soient $f : X \ra S$, $g : Y \ra S$ des applications continues
entre espaces topologiques. Alors le topos $\widetilde{X \times_S Y} $ des faisceaux sur le produit
fibré usuel $X \times_S Y$ représente le produit fibré $\gtilde X \times_{\gtilde S} \gtilde Y$ : le
morphisme naturel $\widetilde{X \times_S Y} \ra \gtilde X \times_{\gtilde S} \gtilde Y$ est une
équivalence. On définit en effet un quasi-inverse en observant que les ouverts de $X \times_S Y$ de
la forme $U \times_V W$ forment une base. 

\item \textit{Schémas}. Soient $f : X \ra S$, $g : Y \ra S$ des morphismes de schémas. Désignons par
l'indice $\zar$ (resp. $\et$) le topos zariskien resp. étale) associé. Du fait de (1),
le morphisme naturel $(X \times_S Y)_{\zar} \ra X_{\zar} \times_{S_{\zar}} Y_{\zar}$ n'est pas une
équivalence en général. De même, le morphisme naturel $(X \times_S Y)_{\et} \ra
X_{\et} \times_{S_{\et}} Y_{\et}$ n'est pas une équivalence en général, même
si $X$, $Y$, $S$ sont les spectres de corps : si $S = \SP k$, $S_{\et}$ est équivalent au
topos classifiant $BG$ du groupe profini $G = \mathrm{Gal}(\overline{k}/k)$, où $\overline{k}$ est
une clôture séparable de $k$, et la formation de $BG$ commute aux produits fibrés.  \section{Topos
évanescents et co-évanescents}
\end{enumerate}
\end{exemples}


\subsection{}\label{11-déf. topos évanescents} Soient $f : X \ra S$,
$g : Y \ra S$ des morphismes de topos comme en \ref{11-site}. Le produit orienté
\begin{equation}
X \ori_S S, \label{11-évanescent}
\end{equation}
où $S \ra S$ est le morphisme identique, s'appelle le
\textit{topos évanescent} de $f$. Il est étudié dans \cite{AG_Tokyo-Kyoto@Laumon} et
\cite{Cycles@Orgogozo}. Le produit orienté
\begin{equation}
S \ori_S Y, \label{11-co-évanescent}
\end{equation}
où $S \ra S$ est le morphisme identique, joue un rôle important dans les travaux de
Faltings sur les théorèmes de comparaison $p$-adiques et la correspondance de Simpson $p$-adique
(\textit{topos de Faltings}) (cf. \cite{almost-etale@Faltings}, \cite{Simpson@Faltings},
\cite{Simpson@Abbes-Gros}, \cite{toposdefaltings@Deligne}, \cite{co-evanescent@AG}). On propose ici de
l'appeler \textit{topos co-évanescent} de $g$. Du topos évanescent de $\Id_S$,
\begin{equation}
\fl S = S \ori_S S, \label{11-S fléché}
\end{equation}
qui est aussi le topos co-évanescent de $\Id_S$, les
produits orientés $X \ori_S Y$ se déduisent par changement de base. Considérons en effet le produit
fibré itéré
\begin{equation}
Z = X \times_S \fl S \times_S Y,
\end{equation} où la flèche de gauche
(resp. droite) de $\fl S$ vers $S$ est $p_1$ (resp. $p_2$). On a des projections naturelles $q_1 : Z
\ra X$, $q_2 : Z \ra Y$ et $m : Z \ra \fl S$, avec des isomorphismes $p_1m \iso fq_1$, $gq_2 \iso
p_2m$. Par composition avec la flèche structurale $\tau : p_2 \ra p_1$ de $\fl S$, on en déduit une
flèche $z : gq_2 \ra fp_1$. D'après\ref{11-existence}, le triplet $(q_1,q_2,z)$ définit donc une
flèche
\begin{equation}
h : Z \ra X \ori_S Y \label{11-h1}
\end{equation}
et des isomorphismes $p_1h
\iso q_1$, $p_2h \iso q_2$, par lesquels $z$ s'identifie à $\tau h$, où $p_1$, $p_2$ désignent les
projections canoniques de $X \ori_S Y$ sur $X$ et $Y$. 

\begin{proposition}
\label{11-produit orienté
et produit fibré} Le morphisme $h$ (\ref{11-h1}) est un isomorphisme. En particulier, il définit des
isomorphismes canoniques
\begin{equation}
X \times_S \fl S \iso X \ori_S S,
\end{equation}
\begin{equation} \fl S \times_S Y \iso S \ori_S Y. 
\end{equation}
\end{proposition}

Il suffit de montrer que $Z$, muni de $(q_1,q_2,z)$ vérifie la
propriété universelle du produit orienté. Soit $T$
un topos muni de morphismes $a : T \ra X$, $b : T \ra Y$, et d'une 2-flèche $t : gb \ra fa$. Par la
propriété universelle de $\fl S$, on en déduit d'abord un unique triplet, formé d'un morphisme $k :
T \ra \fl S$ et d'isomorphismes $p_1k \ra fa$, $p_2k \ra gb$ tels que $t = \tau k$ modulo ces
identifications. Puis, par la propriété universelle des produits fibrés, on en déduit un unique
quadruplet formé d'un morphisme $s : T \ra Z$ et d'isomorphismes $ms \iso k$, $q_1s \iso a$, $q_2s
\iso b$ tel que $zs = t$ modulo ces identifications.  \subsection{}\label{11-psi} Soit $f : X \ra S$
un morphisme de topos. D'après \ref{11-existence}, les morphismes $\Id_X$ et $f$ définissent un
morphisme
\begin{equation}
\Psi : X \ra X \ori_S S \label{11-psi1}
\end{equation}
tel que
\begin{equation}
p_1\Psi = \Id_X, p_2\Psi = f, \tau \Psi = \Id_f, \label{11-psi2}
\end{equation}
où
$\tau : p_2 \ra fp_1$ est la 2-flèche structurale de $X \ori_S S$ :
\begin{equation}
\xymatrix{{} &
X \ar[d]^{\Psi} \ar[dl]_{\Id_X} \ar[dr]^f & {} \\ X \ar[dr]_f & X \ori_S S \ar[l]_{p_1} \ar[r]^{p_2}
& S \ar[dl]^{\Id_S} \\ {} & S & {}} \label{11-psi1bis}
\end{equation}

 Le foncteur $\Psi_*$ s'appelle \textit{foncteur cycles proches}. Pour un objet $(U \ra V \la W)$ du
site $C$ définissant $X \ori_S S$, on a $\Psi^*(U \ra V \la W) = U \times_V W$, où $U \times_V W :=
U \times_{f^*(V)} f^*W$. Si $\Lambda$ est un anneau et $F \in D^+(X,\Lambda)$, le complexe $R\Psi_*F
\in D^+(X \ori_S S,\Lambda)$ (noté aussi $R\Psi F$) s'appelle \textit{complexe des cycles proches}
(de $f$ relatif à $F$). 

L'identité $p_1{}_*\Psi_* = \Id$ définit, par adjonction, un morphisme canonique
\begin{equation}
p_1^* \ra \Psi_*.  \label{11-psi3}
\end{equation}
Pour $F \in D^+(X,\Lambda)$, le cône du morphisme
$p_1^*F \ra R\Psi F$ qui s'en déduit s'appelle le \textit{complexe des cycles évanescents} (de $f$
relatif à $F$) et se note $R\Phi F$. Dans le cas des schémas (munis de la topologie étale), ces
foncteurs, qui généralisent les foncteurs $R\Psi$ et $R\Phi$ de Grothendieck (\sga{7}{i, xiii}),
sont étudiés dans \cite{AG_Tokyo-Kyoto@Laumon} et \cite{Cycles@Orgogozo}. 

Considérons le morphisme de changement de base
\begin{equation}
p_1{}_* \ra \Psi^* \label{11-psi4}
\end{equation} déduit de l'identité $p_1\Psi = \Id_X$, en d'autres termes, le morphisme déduit, par
application de $p_1{}_*$, de la flèche d'adjonction $\Id \ra \Psi_* \Psi^*$, compte tenu de ce que
$p_1{}_*\Psi_* = \Id$. Le résultat suivant est donné sans démonstration dans
\cite{AG_Tokyo-Kyoto@Laumon} :

\begin{proposition}
\label{11-isocan1} Le morphisme \ref{11-psi4} est
un isomorphisme. 
\end{proposition}

On va définir un morphisme
\begin{equation}
\Psi^* \ra p_1{}_*,
\label{11-psi5}
\end{equation}
dont on montrera qu'il est inverse de \ref{11-psi4}. Pour cela, on
définit un morphisme
\begin{equation}
\Id \ra \Psi_*p_1{}_* \label{11-psi6}
\end{equation}
de la
façon suivante. Pour un faisceau $F$ sur $X \ori_S S$ et un objet $Z = (U \ra V \la W)$ du site $C$
de \ref{11-site}, la flèche $F(Z) \ra (\Psi_*p_1{}_*F)(Z)$ est la composée
\begin{equation}
F(Z) \ra
F(U \times_{f^*V} f^*W \ra W \la W) \ra F(U \times_{f^*V} f^*W \ra e_S \la e_S), \label{11-psi7}
\end{equation}
où $W \ra W$ est l'identité, la première flèche est la restriction et la seconde,
l'inverse de l'isomorphisme relatif au recouvrement de type (c) $$ (U \times_{f^*V} f^*W \ra W \la
W) \ra (U \times_{f^*V} f^*W \ra e_S \la e_S).  $$ Le morphisme \ref{11-psi5} est adjoint de
\ref{11-psi6}. Notons $u$ (resp. $v$) le morphisme \ref{11-psi4} (resp. \ref{11-psi5}). On va
montrer que $u$ et $v$ sont inverses l'un de l'autre. L'argument qui suit est dû à Orgogozo. Il
s'agit de montrer que, pour tout faisceau $F$ sur $X \ori_S S$ et tout faisceau $G$ sur $X$, les
applications $\alpha(F,G) = \mathrm{Hom}(u(F),G) : \mathrm{Hom}(\Psi^*F,G) \ra
\mathrm{Hom}(p_1{}_*F,G)$ et $\beta(F,G) = \mathrm{Hom}(v(F),G) :\mathrm{Hom}(p_1{}_*F,G) \ra
\mathrm{Hom}(\Psi^*F,G)$ sont inverses l'une de l'autre.

 L'application $$ \alpha(F,G) : \mathrm{Hom}(\Psi^*F,G) = \mathrm{Hom}(F,\Psi_*G) \ra
\mathrm{Hom}(p_1{}_*F,G) $$ envoie $a : F \ra \Psi_*G$ sur $p_1{}_*a : p_1{}_*F \ra p_1{}_*\Psi_*G =
G$. L'application $a$ est la donnée d'une famille compatible d'applications $a_{(U \ra V \la W)} :
F(U \ra V \la W) \ra G(U \times_V W)$, pour $(U \ra V \la W)$ parcourant les objets de $C$,
``compatible" voulant dire compatible aux flèches de restriction. L'application $\alpha(a)$ est la
famille $a_{(U \ra e_S \la e_S)} : F(U \ra e_S \la e_S) \ra G(U)$, $U$ parcourant les objets de $X$.

L'application $$ \beta(F,G) : \mathrm{Hom}(p_1{}_*F,G) \ra \mathrm{Hom}(\Psi_*F,G) =
\mathrm{Hom}(F,\Psi_*G) $$ envoie $b : p_1{}_*F \ra G$ sur le composé $F \ra \Psi_*p_1{}_*F \ra
\Psi_*G$, où la première flèche est \ref{11-psi6} et et la seconde $\Psi_*b$. L'application
$\beta(b)$ est la famille $\beta(b)_{(U \ra V \la W)} : F(U \ra V \la W) \ra G(U \times_V W)$ est la
composée de \ref{11-psi7} et de $$ (\Psi_*b)(U \times_V W \ra e_S \la e_S) : F(U \times_V W \ra e_S
\la e_S) = (p_1{}_*F)(U \times_V W) \ra G(U \times_V W).  $$

Pour chaque $a : F \ra \Psi_*G$, on a un diagramme commutatif $$ \xymatrix{F(U \ra V \la W) \ar[d]
\ar[r]^{a_{(U \ra V \la W)}} & G(U \times_V W) \\ F(U \times_V W \ra W \la W) \ar[r] & F(U \times_V
W \ra e_S \la e_S) \ar[u]_{a_{(U\times_V W \ra e_S \la e_S)}}}, $$ où la flèche horizontale
inférieure est l'isomorphisme figurant dans \ref{11-psi7}. Dans ce diagramme, le composé des flèches
autres que la flèche horizontale supérieure est $\beta \alpha(a)_{(U \ra V \la W)}$, donc $\beta
\alpha = \Id$. La vérification de $\alpha \beta = \Id$ est triviale également. 

\begin{remarque}\label{11-isocan2} Lorsque les topos $X$ et $S$ sont localement
cohérents, donc en particulier dans le cas des schémas, on peut prouver \ref{11-isocan1} plus
simplement, en se ramenant au cas local. Comme $X$ a assez de points, il suffit de montrer que, pour
tout point $x : pt \ra X$, et tout faisceau $F$ sur $X \ori_S S$, la fibre en $x$ de \ref{11-psi4},
est un isomorphisme. Quitte à remplacer $X$ par son localisé en $x$ (\textit{cf.}
\ref{11-localisation1}), on peut supposer $X$ local de centre $x$. Soit $s : pt \ra S$ l'image de
$x$ par $f$. Soit $F$ un faisceau sur $T = X \ori_S S$. On doit montrer que $$ (p_1{}_*F)_x =
\Gamma(T,F) \ra (\Psi_*F)_x = F_{(x,s)} $$ est un isomorphisme. Soient $S_{(s)}$ le localisé de $S$
en $s$ et $T_{(s)} = T \ori_{S_{(s)}} S_{(s)}$. D'après les résultats de passage à limite invoqués
dans \ref{11-preuve de ch. base}, on a $$ \Gamma(T_{(s)},F) = \mathrm{colim} \Gamma(T_U,F), $$ où
$U$ parcourt les voisinages cohérents de $s$ et $T_U = X \ori_{S_{/U}} S_{/U}$. D'après
\ref{11-localisation4}, les flèches de restriction $\Gamma(T,F) \ra \Gamma(T_U,F)$ sont des
isomorphismes. On peut donc supposer $S$ local de centre $s$. D'après \ref{11-can9}, $T$ est alors
local, de centre $(x,s)$, et $\Gamma(T,F) = F_{(x,s)} = (\Psi^*F)_x$.
\end{remarque}

\subsection{}\label{11-topos de Faltings}
Soit $g : Y \ra S$ un morphisme de topos. Les faisceaux sur le topos co-évanescent $T =
S \ori_S Y$ (\ref{11-co-évanescent} ont une description simple, due à Deligne
\cite{toposdefaltings@Deligne}. Pour un faisceau $F$ sur $T$, la flèche de restriction $F(U \ra V
\la W) \ra F(U \ra U \ra U \times_V W)$ est un isomorphisme. Cela suggère de considérer le site
$C_0$ suivant. La catégorie $C_0$ est celle des flèches $(V \la W)$ au-dessus de $f : Y \la S$, i.
e. des flèches $W \ra g^*V$ de $Y$. On munit $C_0$ de la topologie définie par la prétopologie
engendrée par les familles couvrantes des types (a) et (b) ci-après :

(a) $(V \la W_i)_{i\in I} \ra (V \la W)$, où la famille $(W_i \ra W)_{i \in I}$ est couvrante,

(b) $(V_i \la W_i)_{i \in I} \ra (V \la W)$, où la famille $(V_i \ra V)_{i \in I}$ est couvrante, et
$W_i = V_i \times_V W$.

On montre (\cite[4.10]{co-evanescent@AG}) que $C_0$  est un site de définition de $T$. Un objet $F$
de $S \ori_S Y$, décrit comme un faisceau sur $C_0$, s'interprète comme la donnée d'une famille de
faisceaux $F_V : W \mapsto F(V \la W)$ sur $g^*V$ et de flèches de restriction $F_V \ra
j_{V'V}{}_*F_{V'}$ pour $V' \ra V$, définissant $j_{V'V} : g^*V' \ra g^*V$, satisfaisant la
condition de descente que, pour une famille couvrante $(V_i \ra V)_{i \in I}$, la suite $$ F_V \ra
\prod_{i} j_{V_iV}{}_*F_{V_i} \rightrightarrows \prod_{ii'} j_{V_{ii'}V}{}_*F_{V_{ii'}} $$ soit
exacte, où $V_{ii'} = V_i \times_V V_{i'}$. D'après \ref{11-existence}, les morphismes $g : Y \ra S$
et $\Id_Y$ définissent un morphisme
\begin{equation}
\Psi : Y \ra S \ori_S Y \label{11-copsi1}
\end{equation} tel que

\begin{equation}
p_1\Psi = g , p_2\Psi = \Id_Y , \tau\Psi = \Id_g,
\label{11-copsi2}
\end{equation}
où $\tau : gp_2 \ra p_1$ est la 1-flèche structurale de $S \ori_S Y$ :
\begin{equation}
\xymatrix{{} & Y \ar[d]^{\Psi} \ar[dl]_{g} \ar[dr]^{\Id_Y} & {} \\ S
\ar[dr]_{\Id_S} & S \ori_S Y \ar[l]_{p_1} \ar[r]^{p_2} & Y \ar[dl]^{g} \\ {} & S & {}}.
\label{11-copsi1bis}
\end{equation}
Le foncteur $\Psi_*$, qu'on pourrait appeler \textit{foncteur
cycles co-proches}, se comporte de manière très différente du foncteur cycles proches de
\ref{11-psi1}. En effet, de l'identité $p_2\Psi = \Id_Y$ on déduit, par adjonction, un morphisme
canonique
\begin{equation}
p_2^* \ra \Psi_*, \label{11-copsi3}
\end{equation}
analogue de
\ref{11-psi3}, et l'on a :

\begin{proposition}
\label{11-copsi4} Le morphisme \ref{11-copsi3} est un
isomorphisme. 
\end{proposition}
En particulier, \textit{le foncteur $\Psi_*$ est exact.} Ici, c'est
la flèche de changement de base, déduite de l'identité $p_2\Psi = \Id_Y$,
\begin{equation}
p_2{}_*
\ra \Psi^*, \label{11-copsi5}
\end{equation}
analogue de \ref{11-psi4}, qui n'est pas, en général,
un isomorphisme. On peut donner de \ref{11-copsi4} une démonstration analogue à celle de
\ref{11-isocan1}. Il est plus simple de déduire ce résultat des descriptions explicites suivantes
des foncteurs $p_1^*$, $p_2^*$, $\Psi$ et du morphisme $\tau$. Ces descriptions sont dues à Deligne
\cite{toposdefaltings@Deligne}.  \subsection{}\label{11-descriptions} (a) \textit{Description de
$p_1^*$}. On a $p_1^*V = (V \la g^*V)$ (l'objet $(V \ra e_S \la e_Y)$ de $C$ correspondant à l'objet
$(V \la g^*V)$ de $C_0$). Si $F$ est un faisceau sur $S$, $p_1^*F$ est le faisceau associé au
préfaisceau dont la valeur en $(V \la W)$ est la limite inductive des $F(V')$ suivant la catégorie
des flèches $(V \la W) \ra p_1^*V'$. Cette catégorie ayant $(V \la W) \ra (V \la g^*V)$ pour objet
initial, cette limite est égale à $F(V)$. En d'autres termes, $p_1^*F$ est le faisceau associé au
préfaisceau dont la valeur en $(V \la W)$ est $F(V)$. Dans la description donnée plus haut d'un
faisceau sur $S \ori_S Y$ en termes d'une famille de faisceaux sur les $g^*V$, $p_1^*F$ est la
famille des faisceaux constants $G_V$ sur $g^*V$ de valeur $F(V)$.

\medskip (b) \textit{Description de $p_2^*$.} On a $p_2^*W = (e_S \la W)$. Si $F$ est un faisceau
sur $Y$, $p_2^*F$ est le faisceau associé au préfaisceau dont la valeur en $(V \la W)$ est la limite
inductive des $F(W')$ suivant la catégorie des flèches $(V \la W) \ra p_2^*W'$. Cette catégorie
ayant $(V \la W) \ra (e_S \la W)$ pour objet initial, cette limite est égale à $F(W)$. En d'autres
termes, $p_2^*F$ est la famille des faisceaux $H_V$ sur $g^*V$, où $H_V$ est le faisceau $W \mapsto
F(W)$. 

\medskip (c) \textit{Description de $\tau$.} Si $F$ est un faisceau sur $S$, le morphisme $\tau :
p_1^*F \ra (gp_2)^*F$ est déduit du morphisme de préfaisceaux qui, pour $(V \la W)$ dans $C_0$,
envoie $F(V)$ dans $(gp_2)^*F(V \la W) = (g^*F)(W)$ par le morphisme composé $F(V) \ra (g^*F)(g^*V)
\ra (g^*F)(W)$. On le voit à l'aide de (a) et (b), en explicitant le morphisme $(fp_1)^* \ra
(gp_2)^*$ adjoint du morphisme $\tau$ décrit en \ref{11-tau}.

\medskip (d) \textit{Description de $\Psi_*$.} Si $F$ est un faisceau sur $Y$, et $(U \ra V \la W)$
un objet de $C$, on a $(\Psi_*F)(U \ra V \la W) = F(U \times_V W$. Dans la description de $S \ori_S
Y$ à l'aide du site $C_0$, on a donc $$ (\Psi_*F)(V \la W) = F(W).  $$ Compte tenu de (b), on a donc
$$ \Psi_*F = p_2^*F $$ Cette identification est celle donnée par \ref{11-copsi3}, ce qui prouve
\ref{11-copsi4}.

Notons encore que les points de $T = S \ori_S Y$ sont les flèches $s \la y$, où $s$ (resp. $y$) est
un point de $S$ (resp. $Y$), et que, pour un faisceau $F$ sur $Y$, si $(s \la y)$ est un point de
$T$, on a $$ (p_2^*F)_{(s \la y)} = F_y = (\Psi_*F)_{(s \la y)}.  $$

\begin{remarque}
Comme l'observe Gabber, l'isomorphisme \ref{11-copsi3} implique le théorème de
changement de base \ref{11-ch. base} pour $X = S, f = \Id_S$ sans hypothèse de cohérence sur $S$,
$Y$, et $g$.
\end{remarque}

\ifx\danslelivre\undefined
\bibliography{../fichiers-configuration-latex/bibliographie}
\bibliographystyle{../fichiers-configuration-latex/smfalpha}
\end{document}
\fi



